\documentclass[nynorsk]{scrartcl}

\usepackage{babel}
\usepackage{tikz}

\title[Programmering og Data \\ skisse til EVU-kurs}
\author{Hans Georg Schaathun og Jonas Harang}

\begin{document}

Institutt for IKT og realfag ved NTNU førebur eit nytt etter- og
vidareutdanningskurs i programmering og dataanalyse, herunder maskinlæring.
Kurset er tenkt å gå fyrste gongen i vårsemesteret 2026.

\textbf{Omfang.}  7½ studiepoeng (ECTS) som tilsvarer rundt 200 arbeidatimar.
Dette er fordelt på tre samlingar à 15h og tolv veker à 12h sjølvstudium.
Kurset kan inngå i ein \textbf{Bachelor}-grad.

\paragraph{Målgruppe.}  
Kurset er retta mot folk som er vande med å arbeida med kvantitative data i jobben,
anten det er rekneskap, forretningsanalyser, demografi eller finans.
Elles krev kurset ingen særskilde forkunnskapar.

\paragraph{Læringsutbyte.}  
Etter endt kurs skal studenten kunne henta datasett frå ulike kjelder og bruka
programmeringsverkty for å formattera, analysera, presentera og ikkje minst 
tolka dei innanfor eigen fagdisiplin.

\paragraph{Innhald.}  
Kurset legg vekt på å demonstrera eit rikt monn av ulike teknikkar på ulike
typar data.  Det er i stor grad opp til den einskilde deltakaren å velja kva
teknikkar som er verd 

\paragraph{Vurdering.}  
Grunnlaget for vurdering er ein innlevert mappe der studenten presenterer
utvalde oppgåver som dei har arbeidd med i løpet av semesteret, saman med ein
refleksjon over kva nytteverdi desse oppgåvene har for dei.

\paragraph{Teknisk ramme.}
Hovuddelen av kurset vil bruka Python3 gjennom Jupyter Notebook, saman med
Python-bibliotek som pandas, scikit-learn og torch.  
Me vil òg koma inn på alternativ til Jupyter Notebook (t.d.\ IDE-ar 
--- \emph{Integrated Development Environemt}).
Studentane bruker eigne (bærbare) maskiner, men me gjev vegleiing til installasjon
og oppstart til dei som treng det.



\begin{figure}
   \begin{tikzpicture}
      \tikzstyle{box}=[draw,rounded corners,node distance=2in]
       \tikzstyle{g}=[box,fill=green!20,minimum width=0.8\textwidth]
       \tikzstyle{p}=[box,fill=yellow!20,minimum width=0.7\textwidth,xshift=0.05\textwidth] ;
       \node [g] (g1) {\textbf{Opningssamling} (to dagar).  
                       Koma i gang med verktya.  Overblikk over semesteret.
                       } ;
       \node [p,below of=g1] (p1) {\textbf{Fyrste Periode} (seks veker à 12h).  
           Kvar veke demonstrerer  
       Koma i gang med verktya.  Overblikk over semesteret.} ;
       \node [g,below of=p1] (g2) {\textbf{Midtvegssamling.}  
                          Seminar med erfaringsutveksling frå fyrste periode.
                          Munnleg presentasjon av fyrste del av vurderingsmappa.
                          Overblikk og planleggjing av andre periode.}
       \node [p,below of=g2] (p2) {\textbf{Andre Periode} (seks veker à 12h).
           Koma i gang med verktya.  Overblikk over semesteret.
           } ;
       \node [g,below of=p2] (g3) {\textbf{Oppsummeringssamling.}  \emph{Debrief.}  
            Erfaringsutveksling.  Spørsmål og svar.} ;
   \end{tikzpicture}
   \caption{Kursstruktur.}

\end{figure}

\end{document}
